% The chain connector: a shrouded 2x12 box header, IDC ribbon numbering
% (rfcs/0001-g3/02-connectors.md#chain-connector).
\documentclass[border=6pt]{standalone}
\usepackage{g3-fig}
\begin{document}
\begin{tikzpicture}
  \tikzset{cp/.style={draw, minimum width=1.25cm, minimum height=1.05cm, inner sep=1pt,
                      align=center, font=\ttfamily\scriptsize}}
  % \c{n}{kind colour}{name}: odd pins on the top row, even pins below
  \newcommand\cpin[3]{%
    \pgfmathtruncatemacro\col{(#1-1)/2}
    \pgfmathsetmacro\row{mod(#1,2)==1 ? 0 : -1.15}
    \node[cp, fill=#2] (p#1) at (\col*1.35, \row) {#1\\#3};}

  \cpin{1}{g3Stop}{RUN\_\\FEED}  \cpin{2}{g3Gnd}{GND}
  \cpin{3}{g3Gnd}{GND}           \cpin{4}{g3Ctl}{POS1}
  \cpin{5}{g3Ctl}{POS2}          \cpin{6}{g3Ctl}{POS3}
  \cpin{7}{g3Ctl}{POS4}          \cpin{8}{g3Ctl}{RESET}
  \cpin{9}{g3Ctl}{SYNC}          \cpin{10}{g3Gnd}{GND}
  \cpin{11}{g3Gnd}{GND}          \cpin{12}{g3Stop}{LOOP\_\\OUT}
  \cpin{13}{g3Gnd}{GND}          \cpin{14}{g3Gnd}{GND}
  \cpin{15}{g3Ctl}{INT\_N}       \cpin{16}{g3Gnd}{GND}
  \cpin{17}{g3Bus}{SCL}          \cpin{18}{g3Gnd}{GND}
  \cpin{19}{g3Bus}{SDA}          \cpin{20}{g3Bus}{VIO}
  \cpin{21}{g3Gnd}{GND}          \cpin{22}{g3Gnd}{GND}
  \cpin{23}{g3Gnd}{GND}          \cpin{24}{g3Stop}{LOOP\_\\RET}

  % the two key positions
  \draw[very thick, dashed] (p16.south west) rectangle (p16.north east);
  \draw[very thick, dotted] (p21.south west) rectangle (p21.north east);
  \node[note, anchor=north, align=center] at (p16.south) {IN: no post};
  \node[note, anchor=south, align=center] at (p21.north) {OUT: no post};

  % shroud with its polarising notch on the odd-pin side
  \draw[thick] ($(p1.north west)+(-0.25,0.65)$) rectangle ($(p24.south east)+(0.25,-0.55)$);
  \fill[white] ($(p11.north west)+(0,0.66)$) rectangle ($(p13.north east)+(0,0.62)$);
  \draw[thick] ($(p11.north west)+(0,0.65)$) -- ++(0,-0.25) -| ($(p13.north east)+(0,0.65)$);
  \node[note, anchor=south] at ($(p11.north)!0.5!(p13.north)+(0,0.7)$) {polarising notch};

  % key
  \begin{scope}[shift={(0, -3.1)}]
    \foreach \x/\c/\t in {0/g3Stop/12\,V stop circuit, 3.1/g3Gnd/ground,
                          5.3/g3Ctl/control and position, 9.0/g3Bus/I\textsuperscript{2}C and VIO}
      { \node[key, fill=\c, anchor=west] at (\x, 0) {};
        \node[note, anchor=west] at (\x+0.4, 0) {\t}; }
    \node[note, anchor=north west, align=left, text width=15.5cm] at (0, -0.45)
      {Pin \emph{n} faces \emph{n}$\pm$1 across the rows and \emph{n}$\pm$2 along its row,
       so each 12\,V pin (1, 12, 24) has ground on every neighbour.
       A ribbon blocks hole 21 at its upstream end and hole 16 at its downstream end,
       so it mates only OUT to IN; the terminator blocks hole 21 and fits only an OUT.};
  \end{scope}
\end{tikzpicture}
\end{document}
